\section*{Problem 8}

Suppose \(n = pqr\) is the product of three distinct primes. How would an RSA-type
scheme work in this case? What relation would \(e\) and \(d\) satisfy?

\textbf{Solution.}

If \(n = pqr\) with distinct primes \(p, q, r\), then an RSA-type system works the same way as usual,
but the Euler totient becomes
\[
\varphi(n) = (p-1)(q-1)(r-1).
\]

To choose an encryption exponent \(e\), we require
\[
\gcd(e, \varphi(n)) = 1.
\]

Then the decryption exponent \(d\) must satisfy
\[
ed \equiv 1 \pmod{\varphi(n)}.
\]

Encryption and decryption work as usual:
\[
c \equiv m^e \pmod{n}, \qquad m \equiv c^d \pmod{n}.
\]

The only change from standard RSA (two primes) is the definition of \(\varphi(n)\).  
Otherwise, the construction and correctness of the scheme remain the same.

